% !TeX program = xelatex
% !TeX encoding = UTF-8
\documentclass[runningheads,orivec]{llncs}
% 必须在加载中文字体前选用 XeLaTeX，避免误用 pdfLaTeX 后查找 unihei 字体。
\RequirePackage{iftex}
\RequireXeTeX
\usepackage[UTF8,scheme=plain,fontset=fandol]{ctex}
\usepackage{amsmath,amssymb,graphicx,booktabs,tabularx,array,longtable}
\usepackage{algorithm,algpseudocode}
\usepackage{tikz}
\usetikzlibrary{arrows.meta,positioning,fit}
\usepackage{url}
% 本地TeX发行版未安装xurl；允许长来源路径在字母和数字间换行。
\makeatletter
\g@addto@macro\UrlBreaks{\do\a\do\b\do\c\do\d\do\e\do\f\do\g\do\h\do\i\do\j\do\k\do\l\do\m\do\n\do\o\do\p\do\q\do\r\do\s\do\t\do\u\do\v\do\w\do\x\do\y\do\z\do\A\do\B\do\C\do\D\do\E\do\F\do\G\do\H\do\I\do\J\do\K\do\L\do\M\do\N\do\O\do\P\do\Q\do\R\do\S\do\T\do\U\do\V\do\W\do\X\do\Y\do\Z\do\0\do\1\do\2\do\3\do\4\do\5\do\6\do\7\do\8\do\9}
\makeatother
\usepackage[hidelinks,unicode,bookmarksdepth=2]{hyperref}
\usepackage{fvextra}
\floatname{algorithm}{算法}
\renewcommand{\figurename}{图}
\renewcommand{\tablename}{表}
\renewcommand{\refname}{参考文献}
\renewcommand{\abstractname}{摘要：}
\renewcommand{\keywordname}{{\bfseries 关键词：}}
\renewcommand{\contentsname}{目录}
\algrenewcommand\algorithmicrequire{\textbf{输入：}}
\algrenewcommand\algorithmicensure{\textbf{输出：}}
\newcommand{\clip}{\operatorname{clip}_{[0,1]}}
\newcommand{\AQ}{\operatorname{AQ}}
\newcommand{\gap}{\operatorname{gap}}
\newcommand{\FE}{\mathrm{FE}}
\newcommand{\JS}{\textsc{GP-J}}
\newcommand{\CT}{\textsc{GP-C}}
\newcommand{\FACO}{\textsc{FACO-N}}
\newcommand{\Ueight}{\textsc{GPU-U8}}
\newcommand{\Usixteen}{\textsc{GPU-U16}}
\newcommand{\Rsixteen}{\textsc{GPU-R16}}
\newcolumntype{L}[1]{>{\raggedright\arraybackslash}p{#1}}
\newcolumntype{Y}{>{\raggedright\arraybackslash}X}
\setlength{\emergencystretch}{1em}
\newcommand{\gentable}[1]{\input{generated/#1.tex}}
\newcommand{\genfigure}[1]{\resizebox{\textwidth}{!}{\input{generated/#1.pgf}}}
\DefineVerbatimEnvironment{TreeCode}{Verbatim}{fontsize=\small,breaklines=true,breakanywhere=true}

\begin{document}
\title{补充材料：GP学习FACO搜索控制策略}
\titlerunning{GP--FACO实验与规则补充材料}
\author{匿名作者}\authorrunning{匿名作者}\institute{}\maketitle
\begin{abstract}
本材料补充主稿的实现对应、预实验记录、六个选定GP个体的完整树结构、真实输入输出、逐seed结果和计算成本，并整理此前完成的实验。正式比较使用两种表示各3个GP seed、128个individual、50代训练及冻结测试。实验沿革单独记录旧版32蚂蚁实验、暂停轮次和工程加速；各阶段的设置与结果分别呈现。实验结果截至2026年9月10日。
\end{abstract}
\setcounter{tocdepth}{1}
\section*{目录}
\makeatletter
\begingroup
\let\l@title\@gobbletwo
\let\l@author\@gobbletwo
\let\authcount\@gobble
\@starttoc{toc}
\endgroup
\makeatother
\clearpage

\section{实验沿革与证据范围}
\begin{longtable}{@{}L{20mm}L{47mm}L{49mm}@{}}
\caption{已开展工作的分类。完成工程检查、开始训练和完成科学比较是不同状态。}\label{tab:catalog}\\
\toprule 阶段 & 做了什么 & 状态与可支持的结论\\ \midrule\endfirsthead
\toprule 阶段 & 做了什么 & 状态与可支持的结论\\ \midrule\endhead
早期底座与候选图 & 数据读取、CPU/CUDA路径构造、状态分叉、候选图匹配、Hard/Escape机制、计数预算和任务恢复。 & 工程记录位于results/v1；证明相应功能曾被检查，不证明E1--E4的算法假设。\\
旧32蚂蚁实验 & E1反馈对照和E3候选图实验曾启动GP；E3的四组静态调参完成。 & 旧GP未完成50代及最终验证，测试未开放。预算、算子和硬件不同，不与新结果拼接。\\
2022 FACO复现 & 六个TSPLIB实例，各30种子，原始整数距离，5000迭代。 & 180次运行完成。用于核查原生算法与论文参数；不是合成测试集上的GP收益。\\
连续距离与精度 & 原生FACO连续距离适配；exact、fp32、fp32\_fast测速门槛。 & 两个FP32搜索候选均未满足加速门槛，保留exact。GP运算仍为FP32。\\
迭代预算预实验 & 两个规模、24种人工设计策略、32个预实验实例、5个ACO种子，比较预算与排名。 & 选择训练H1000，验证／测试H5000。预实验策略不是后来进化出的GP个体。\\
种群并行与旧GP & 个体、实例、蚂蚁展开并行，预实验任务分片；旧单树运行出现浅树和无效交叉。 & 同卡加速证据与保存程序回放完整；旧训练约13代时暂停，保留检查点。\\
表示10代预实验 & 联合单树／条件三树各3种子，每次128个体、10代；预实验实例集终点评价。 & 六次完成，218任务无失败。仅探索性预实验结果，无正式测试。\\
GPU布局预实验 & 同卡四组种群样本、2/4/8 warp比较、六GPU分片大小与完整H1000验证。 & 选2 warp和64个体分片；同卡输出和工作计数保持。\\
表示50代正式轮次 & 两种表示各3种子，从头训练；轻量监控、两级选模、全部冻结后测试。 & 六次训练完成，TSP500及TSP1000测试完成，2382任务无失败、无重试。主稿主要证据。\\
冻结结果诊断 & 64训练情境及144条真实决策的CUDA评分回放、反馈0/1干预、重启累计量分析。 & 无新增ACO FE。证明部分动作对反馈敏感；未完成完整求解的反馈消融。\\
\bottomrule
\end{longtable}

旧E3静态终态记录记载5280次调用、692\,060\,160 FE和168\,960条输出路径，四组为alpha/POPMUSIC候选与Hard/Escape组合，无失败求解。这里输出路径数与FE是不同统计量。该记录属于32蚂蚁旧协议，只有预实验调参结果，不能与本稿128蚂蚁、软候选后备访问的结果比较。旧历史报告还记载E1的10次GP及E3的20次GP已启动；其保存前缀不构成完整实验结论。

主稿未完成的科学问题包括：反馈的完整求解消融、区域／重启／强度的独立贡献、候选先验与逃逸机制的正式因子比较，以及更多规模和分布上的迁移。现有E1--E4设计文档保留这些研究范围，当前50代表示比较不能代替它们。

\section{参数明细与实现语义}
\subsection{阅读符号与代码字段}
主稿用短符号展示tree，本文用相同符号绘制所有选定GP个体。表\ref{tab:symbols}给出与保存记录的对应，方便将数学表达式、真实输入和程序连接起来。图中内部节点为function，浅色叶子为terminal；ERC用$c_i$标记，并在树下列出完整FP32数值。
\begin{table}[htbp]\caption{论文符号与冻结输入字段。}\label{tab:symbols}
\centering\small
\begin{tabular}{@{}clcl@{}}
\toprule 符号 & 字段 & 符号 & 字段\\ \midrule
$p$ & \texttt{progress} & $s$ & \texttt{stagnation}\\
$\rho$ & \texttt{return\_rate} & $w$ & \texttt{ls\_work}\\
$r$ & \texttt{restart} & $u$ & \texttt{mne\_level}\\
$\delta_q$ & \texttt{ref\_gap} & $d_q$ & \texttt{ref\_diff}\\
$e_R$ & \texttt{region\_excess} & $a_R$ & \texttt{archive\_disagreement}\\
$\tau_R$ & \texttt{pheromone\_strength} & $d_R$ & \texttt{region\_dispersion}\\
\bottomrule
\end{tabular}
\end{table}

\subsection{实现对应与原生FACO适配}
GP的terminal计算、function定义、初始化、遗传算子、参数及评价流程已完整列于主稿。上表保留论文符号与保存字段的对应，供阅读原始GP表达式和输入输出记录时使用；后文的来源索引列出参数清单及实现入口。

\FACO{}在合成数据上使用连续FP64欧氏距离，将依赖距离取整的KD-tree邻域实现替换为近邻列表；恒等2-opt移动赋零增益，边集合不变的3-opt移动被排除，以避免浮点误差触发虚假改善循环。该适配保留原生24条初始化路径及2-opt/3-opt初始化。GPU底座则采用2024后续实现的节点重定位构造及单路径初始化。原始整数距离的TSPLIB复现单独报告，三者的路径质量与耗时均按各自设置解释。

GPU构造使用权重$\tau_{ij}d_{ij}^{-\beta}$在主要候选中选择节点；主要候选不可用时，使用随后64个近邻构成的备用表及后备选择。候选表规定优先搜索范围。备选参考路径的边差异在64个节点上采样，优先选差异最大的合格archive路径；随后以路径质量和固定次序处理并列。区域0和1各自随机选择起点，再沿参考路径取连续片段；区域2从随机节点沿主要近邻图扩展。区域在动作执行前生成，动作执行后才重置对应反馈。

\section{原生FACO复现与预实验}
\subsection{2022 FACO的TSPLIB复现}
采用TSPLIB\cite{reinelt1991}中六个实例，保持原始整数距离与2022原生实现\cite{skinderowicz2022}，各运行30个随机种子（1--30），5000迭代，CPU 8线程，24条初始路线。表\ref{tab:tsplib}的差距针对这些TSPLIB实例的已知最优值；它与合成数据的“参考差距”定义依据不同。
\begin{table}[htbp]\caption{原生FACO的180次TSPLIB运行。时间为现有运行记录，包含初始化。}\label{tab:tsplib}
\centering\small\gentable{tsplib}\end{table}
这组结果支持原生复现及规模参数的使用。硬件、随机过程和运行环境未与原论文逐项等同，故不要求逐次路径或耗时复现，也不据此声称GPU底座与原始FACO逐操作等价。

\subsection{迭代预算预实验的完整结果}
主稿的迭代预算预实验使用24种人工设计搜索策略，在每个规模32个预实验实例和5个ACO种子上产生求解曲线。改善保留率定义为
\begin{equation}
 R(H)=\frac{\bar g(0)-\bar g(H)}{\bar g(0)-\bar g(5000)},
\end{equation}
并计算各策略在$H$与5000下的平均差距排名相关。需两个规模均满足$R\ge0.95$且Spearman相关$\ge0.9$，取最小候选预算。表\ref{tab:horizon}给出全部候选。TSP500初始化平均差距5.83718\%，H5000为0.19343\%；TSP1000分别为5.85762\%和0.23266\%。
\begin{table}[htbp]\caption{人工设计搜索策略的迭代预算预实验。保留率衡量相对初始化的改善，不是终点差距的保留比例。}\label{tab:horizon}
\centering\small\gentable{horizon}\end{table}
H500虽然保留较多改善，但排名相关未达标；H1000在两个规模均通过。最终仅训练TSP500。对不依赖进度的人工设计策略，可以从长运行检查点取得短预算结果；进化GP使用$(t-1)/H$，其长运行前缀不必等价于独立短预算求解，后者需重新评价。该限制不能由本次预实验消除。

\subsection{数值精度预实验}
表\ref{tab:precision}比较两种低精度搜索候选与exact后端。四个样本分别为n500/n1000的固定MNE8区域3及代表性GP；静态区域3只在16个随机节点中取起点，不等同于正式\Ueight{}的全节点均匀起点。候选需先达到至少1.10的加速门槛，再进入完整质量门槛。两种候选均未通过速度条件，因此没有足够依据更换exact；不能把这一步写为已经发现FP32质量劣化。
\begin{table}[htbp]\caption{数值精度预实验的运行时间比较，加速比为exact耗时／候选耗时。}\label{tab:precision}
\centering\small\gentable{precision}\end{table}

\section{工程加速、旧GP问题与10代预实验}
\subsection{已有加速证据及比较边界}
早期同卡TSP500对比采用128个体、16实例、1个ACO种子：H100从22.076秒降至12.377秒，H1000从252.661秒降至137.522秒，保存轨迹相同。该结果来自v2构建，涉及档案／反馈更新和构造搜索等热点优化。后续v3以修复表示后的初代和第10代种群为样本，结果见表\ref{tab:speed}。四种样本采用相同A5000，在预热后重复3次取中位数；总预算均为26\,214\,400 FE。
\begin{table}[htbp]\caption{最新v3同卡布局优化。g10为10代预实验末代样本，不是50代末代种群。}\label{tab:speed}
\centering\small\gentable{speed}\end{table}
每block 2、4、8个warp的四种样本比较均支持2 warp。六GPU、H100的分片16、32、64、128个体墙钟中位数依次为60.590、49.780、44.305、48.280秒，选择64。六个真实第10代种群的完整H1000验证共1\,572\,864\,000 FE，六GPU墙钟468.877秒、GPU任务时间和2375.953秒；这是含worker启动的一次完整预算测量。旧版加速比不能与新版百分比连乘为已测得的端到端加速比。

RMTGP-ACO的展开思路为本项目的并行设计提供参考，但现有历史比较并非同卡同工作量。已有报告中，RMTGP的TSP500 AS、ACS、MMAS三种算法平均训练总时间为14.83、2.92、2.63小时，其硬件为RTX PRO 5000 Blackwell、搜索后端为fp32\_fast；本稿为A5000与FP64搜索，评价预算、初始化和终点评分也不同。故本文不将这些数字作为“快于RMTGP”的证据。具体历史记录及配置入口见来源索引。

\subsection{旧单树繁殖问题}
旧v2单树训练出现浅树与交叉无变化。保存的父代、随机状态和后代程序允许直接回放代际变化；表\ref{tab:old}来自已经完成的回放，不新增求解。第9到10代，三个种子交叉均完全没有改变程序，124个后代中只有21--29个新程序。第12到13代仍有相同现象。该证据解释了为何需要修改初始化、允许根交叉、偏好非终端、改善叶子变异并控制行为重复。
\begin{table}[htbp]\caption{旧v2代际程序回放。所有保存的下一代程序均匹配回放结果。}\label{tab:old}
\centering\small\gentable{old_variation}\end{table}
该诊断与修复后的多样性改善之间存在版本差异；没有只改一个算子的完整对照，因而不能估计每条修复的独立效果。旧checkpoint保留且未向新轮次续接适应度。

\subsection{修复表示后的10代预实验}
预实验与正式轮次使用相同的两种表示、3个GP种子、128个体和每代16实例$\times$1个ACO种子$\times$1000迭代，但只运行10代。在固定32个预实验实例$\times$3个ACO种子$\times$5000迭代上评价第10代训练最优个体。表\ref{tab:pilot}列出六次结果。
\begin{table}[htbp]\caption{10代预实验，仅使用预实验实例集，不是正式测试结果。}\label{tab:pilot}
\centering\small\gentable{pilot}\end{table}
两种表示的预实验差距均值分别为0.13172\%与0.16389\%；同一评价集\FACO{}、\Ueight{}、\Rsixteen{}、\Usixteen{}分别为0.20047\%、0.14446\%、0.14239\%、0.15469\%。预实验显示修复后仍有98--112种动作向量，值得完成更长训练；它没有证明单树在独立测试中稳定占优。后来的50代选模验证集和正式测试集不同，不能用两张表的均值直接计算“多训练40代”的收益。

\section{正式轮次的逐运行结果}
\subsection{选中个体与训练多样性}
表\ref{tab:selected}中的个体编号包含首次产生的代数和种群位置。“快排”是该个体在本次运行快速验证中的名次，完整验证使用不同64实例、3个ACO种子与H5000。每个测试均值包含1280次求解，未只挑选表现最好的GP种子。
\begin{table}[htbp]\caption{六个冻结个体。验证和测试列均为参考差距（\%）。}\label{tab:selected}
\centering\small\gentable{selected}\end{table}
\begin{table}[htbp]\caption{第50代种群多样性。程序、行为和适应度列为不同值的数量；高度是每个个体最大树高的中位数。}\label{tab:diversity}
\centering\small\gentable{diversity}\end{table}
\begin{figure}[htbp]\centering\genfigure{training_absolute}
\caption{六次运行的原始训练差距。训练最优个体和中位数都随每代评价集改变，不能把波动直接解释为遗传退化；本图补充主稿中的同一评价集对照差值。}\end{figure}
\begin{figure}[htbp]\centering\genfigure{monitor}
\caption{固定监控验证集上的在线monitoring。每5代评价一次训练最优个体，16实例、ACO seed 17、H1000；细线为三个GP seed，粗线为均值。它与训练后执行的quick validation具有不同的实例集和评价频率。}\end{figure}

\subsection{完整冻结树表达式}
下列树图直接由存档后缀IR生成，保留全部function、terminal和拓扑，没有代数化简。符号见表\ref{tab:symbols}，运算定义、FP32常量、裁剪与并列规则与主稿一致。\CT{}的seed 2207中，区域树只读取$a_R$，直接选择archive disagreement较高的区域；该树与另外两棵树共同构成27节点individual。单个角色的短表达式与整个种群的结构多样性是不同观察对象。
\gentable{trees}

\section{真实决策与反馈诊断}
联合树对完整动作评分，条件三树按角色输出分数向量；最终输出都是离散动作而非新路径。下面给出各一种已保存的真实决策，保留足够小数显示近并列。未显示的候选输入仍可从\path{results/v3/representation-50gen/decision_examples.json}读取，包含12$\times$32特征矩阵、合法性掩码及全评分或分阶段评分。
\gentable{examples}

\subsection{全部反馈置零／置一结果}
表\ref{tab:diag}中每格为“置零改变动作数／置一改变动作数”。在训练64情境、TSP500回放12情境和TSP1000回放12情境上分别计算。训练情境没有对应的历史实际执行动作，因此其保存动作匹配标志不适用；真实回放的144个动作全部匹配生产CUDA评分器。全部528个原始动作取MNE16，所有干预均未改变MNE。
\begin{table}[htbp]\caption{冻结状态反馈干预，不重新运行ACO。}\label{tab:diag}
\centering\small\gentable{feedback_full}\end{table}
一些特征出现在表达式中，却可能只对候选分数添加公共偏置，或因数值并列而不改变选项。因此，表达式出现频率、动作敏感性和质量因果效应需要分别报告。\CT{}的seed 3313参考路径树在实数下等价于$ws^2\delta_q r/\sqrt2$。保留路径时$r=0$，得分为零；当备选路径的$\delta_q$与反馈$s,w$使切换得分超过并列阈值时，规则选择$r=1$。主稿真实示例中两项得分均为零，按候选编号选择保留路径。

\subsection{完整求解的累计重启}
表\ref{tab:restart500}--\ref{tab:restart1000}使用测试求解的完整累计状态，而不是12情境采样的重启比例。每行1280次求解，每次5000迭代。“零重启次数”是完全没有重启的求解数。均值和中位数的差异表明部分GP策略具有长尾重启频率。
\begin{table}[htbp]\caption{TSP500累计重启次数。}\label{tab:restart500}
\centering\small\gentable{restarts500}\end{table}
\begin{table}[htbp]\caption{TSP1000累计重启次数。}\label{tab:restart1000}
\centering\small\gentable{restarts1000}\end{table}
六个GP策略在两个规模的终点全局停滞计数第10百分位均超过32，说明大量终点的归一化停滞输入已饱和。这不意味着完整轨迹每步都饱和，也不能仅凭终点状态判断控制失败。

\section{完整统计与资源成本}
\subsection{配对统计补充}
表\ref{tab:primary}列出TSP500的9项改善、区间与校正$p$值，对应主稿的区间图。表\ref{tab:wins500}补充胜平负及逐求解相对路径改善；表\ref{tab:transfer}和表\ref{tab:wins1000}列出TSP1000的全部描述性比较。胜平负在每实例平均GP/ACO seed后计算，不以1280个求解视为独立样本。相对路径改善是$100(C_b-C_a)/C_b$的平均值，与平均gap的比例不同。
\begin{table}[htbp]\caption{TSP500完整配对比较，$\Delta$及区间单位为百分点；$p_H$在9项主检验中共同Holm校正。}\label{tab:primary}
\centering\small\gentable{primary}\end{table}
\begin{table}[htbp]\caption{TSP500胜平负与相对路径长度改善。}\label{tab:wins500}
\centering\small\gentable{wins500}\end{table}
\begin{table}[htbp]\caption{TSP1000完整描述性比较，单位为百分点；不进行额外主检验，$p_H$留空。}\label{tab:transfer}
\centering\small\gentable{transfer}\end{table}
\begin{table}[htbp]\caption{TSP1000胜平负与相对路径长度改善。}\label{tab:wins1000}
\centering\small\gentable{wins1000}\end{table}

\subsection{训练、监控、选模与测试工作量}
表\ref{tab:costs}前三个数值列均为GPU任务时间求和，单位分钟；墙钟为训练起止间隔，包含并行重叠与等待，不能与GPU时间相加。“代评价等待”是协调器等待该代评价收齐的均值，不包含一代所有管理步骤。监控每次只评训练最优个体，快速验证最多177个去重候选，完整验证只评4个，因此其预算远小于训练。按候选上限估算，监控加选模FE不超过训练的9.4375\%；实际时间比例略高，反映候选的搜索工作量和阶段开销不同。
\begin{table}[htbp]\caption{正式轮次每次运行的成本。训练、监控、选模单位为GPU分钟。}\label{tab:costs}
\centering\small\gentable{costs}\end{table}
\begin{table}[htbp]\caption{TSP500测试吞吐及每FE搜索工作。总摊销秒为1280次求解的批处理时间分摊之和。}\label{tab:testcost500}
\centering\small\gentable{testcost500}\end{table}
\begin{table}[htbp]\caption{TSP1000测试吞吐及每FE搜索工作，口径同上。}\label{tab:testcost1000}
\centering\small\gentable{testcost1000}\end{table}
表中CPU原生对照未混入GPU时间比较；其1280次求解累计时间分别为4750.44和6397.49秒。GPU任务分布于不同服务器A5000，最多32个实例--种子对批处理，摊销时间可用于描述本次工作流吞吐，不能当成单实例独占设备的求解延迟。最新工程验证记录包括20项CTest、303项Python测试通过及8项跳过；这是构建验证范围，跳过项不计为通过，也不代替科学对照。

\section{数据与实现来源索引}
本材料的表格与曲线由\path{scripts/build_evostar_report.py}生成，主要读取已经版本化的精简JSON。完整日志、tour、数据库和checkpoint仍保存在artifacts中，不复制到论文目录。历史数值后端和迭代预算预实验的精简快照保存在\path{results/v3/evostar/evidence.json}，记录原始来源路径；默认编译不依赖那些未版本化的历史大文件。

\begin{longtable}{@{}L{35mm}L{81mm}@{}}
\caption{论文结论与来源的对应关系。路径均相对GPLSACO根目录。}\\
\toprule 内容 & 来源\\ \midrule\endfirsthead
\toprule 内容 & 来源\\ \midrule\endhead
正式配置与划分 & \path{results/v3/representation-50gen/campaign.json}；同目录\path{panels.json}；\path{docs/experiments/representation_50gen_v3.md}。\\
训练曲线、选中IR、统计 & \path{results/v3/representation-50gen/summary.json}。\\
逐运行分析、资源与重启 & \path{results/v3/representation-50gen/analysis.json}。\\
输入输出与反馈干预 & 同目录\path{decision_examples.json}及\path{selected_feedback_diagnostic.json}。\\
原生180次复现 & \path{results/v2/faco_2022_tsplib.json}。\\
迭代预算与数值精度预实验 & \path{results/v3/evostar/evidence.json}；原始路径在该文件的source\_paths中。\\
最新工程速度 & \path{results/v3/representation_performance.json}及\path{representation_full_verification.json}。\\
10代预实验 & \path{results/v3/pilot_analysis.json}；\path{results/v3/representation-pilot/}。\\
旧GP回放与测速 & \path{results/v3/historical_variation_replay.json}；\path{results/v2/tsp500_performance.json}。\\
旧32蚂蚁阶段 & \path{docs/archive/v1/README.md}；该目录reports；\path{results/v1/e3_static_terminal_results.json}。\\
RMTGP历史对照 & \path{docs/reports/}中的《训练速度与信号分析.md》及其指向的只读MTGP\_ACO运行记录；硬件和工作量不匹配。\\
GP与条件决策实现 & \path{python/gp_faco/diverse_evolution.py}、\path{python/gp_faco/controller.py}、\path{cuda/gp_score.cuh}。\\
构造与状态实现 & \path{cuda/faco_device.cuh}、\path{cuda/control_state.cuh}、\path{cuda/colony_state.cuh}及相关batch实现。\\
统计实现 & \path{python/gp_faco/representation_statistics.py}；方法间同步重采样，九项主检验。\\
\bottomrule
\end{longtable}

正式冻结方法保持不变；如果后续实施新的消融或参数修改，应形成新实验版本，不能回写为本轮测试前已确定的设计。
\bibliographystyle{splncs04}\bibliography{references}
\end{document}
